\chapter{Conclusion}
\label{chap:conclusion}
\section{Summary}
\label{sec:summary}
Many open questions and insightful measurements can be performed
using dileptonic decays of top-quarks. While the channel has a low
branching ratio, it can over a particularly clean event environments.
The recent analyses of quantum effects and quasi-bound state effects
near the production threshold, have motivated further investigations
into this region. The dileptonic decays provide unique and
unparalleled precision for tests of spin-structures in the involved
physics-processes.
Despite offering a clean environment for analyses, the channel has it
own challenges. The presence of two neutrinos from the leptonic decay
of the $W$-bosons leaves $6$ unknown momenta components. Furthermore,
selecting and assigning the $b$-jets to their parent top-quarks
provides additional ambiguities. While the neutrino regression has
been studied in the past, the assignment of the neutrinos is not only
understudied, but the two tasks are also inherently intertwined. Many
variables that are essential for the measurements of the internal
structure, require a full reconstruction of the top-quarks four
momenta. Thus, the reconstruction and its performance is an essential
aspect for the sensitive of many studies of the top-quark pairs.

To address the jet-assignment, a transformer-based machine-learning
model was developed and trained. Different architectures and layouts
have been tested, but showed very similar performance -- despite
incorporating distinct symmetries of the kinematic objects. The
developed model was found to provide substantial improvement in the
efficiency of the jet-assignment. This increased accuracy propagates
to smaller errors for the reconstruction, and lastly also improved
resolutions for these variables, which ultimately increases the
sensitivity of measurements.
Apart from the general improvements that were observed, a more
detailed investigation revealed which regions of phase-space benefit
the most from the reconstruction using the ml-approach.
Notably, the region near the production-threshold proofed to be
particularly challenging. For this region, no substantial improvement
in accuracy compared to the baseline-methods was observed. Some
kinematic features, making reconstruction in this regime difficult
could be identified, but a share of these effects remains elusive.

To further improve the event reconstruction, and resolve the
interdependencies of the two tasks, a model capable to provide
jet-assignment and neutrino regression was developed. The underlying
model architecture extended the previous layout by an additional
output head for the neutrino momenta. While one initial aim was to
resolve the interdependency, the improvements by a combined approach
were limited. Comparing the performance of models performing only
either of the two tasks with the multi-task model, showed no
improvement for accuracy of jet-assignment and only minor
improvements for the regression.
It was, however, found that some cross-talk between the tasks exist.
Therefore, already existing tools that provide neutrino-regression,
both analytical and ml-based, can be expected to provide valuable
information on the jet-assignment.
Further, the developed methods could improve both the error of the
neutrino regression, and the resolution of derived variables -- such
as spin-sensitive markers, and the invariant mass of the
\ttbar-system. However, despite the improvement in the variable
reconstruction, the model was found to come with strong biases in the
reconstructed variables, which led to significant distortions in the
distributions of reconstructed variables. Some measures to correct
such distortions could be identified and discussed. Yet, the
underlying cause for these effects was found to be, most likely, the
probabilistic nature of the neutrino momenta. Since the neutrinos are
inherently underdetermined, they are most accurately modelled by a
conditional probability density. Training a regression model
collapses the output to the mean of this probability distribution,
which can lead to biases in the computation of non-linear variables.
This can be mitigated by modelling and propagating the conditional
density from the neutrino momenta to reconstructed variables and
averaging over their distributions instead. Generally, this would
provide a more complete representation of the kinematics.

Lastly, the performance of the methods developed here was showcased
by reproducing the analysis of bound-state effects at the production
threshold in a highly simplified setting. An Asimov-fit, aimed to
assess the potential impact on an analysis, found that the
uncertainties associated with the signal strength are significantly
decreased by methods with enhance resolution such as $\nu^2$-Flows
and the MultiModalTransformer, as it was presented in this work. The
latter is able to achieve a 3-fold increase in the significance of
the signal strength, which presents are remarkable improvement. In
practice, however, this will be limited by contributions from
systematic uncertainties. Moreover, the enhanced reconstruction of
the MultiModalTransformer comes with a stark limitation of its own:
reconstructed variables of the top-quarks such as the invariant mass
of the \ttbar-system and the spin-sensitive \chel and \chan comes
with significant biases and distortions in the shape. While it was
shown that this does not impact -- as data and simulation are
distorted in the same way, this can be critical for analyses that
extract information from the distributions directly. One notable
example is the measurement of quantum effects in Reference
\cite{atlas_entanglement}, which used the slope of the
\chel-distribution as its marker for entanglement. Hence, additional
consideration would be necessary for such an analysis. Especially the
distributions of reconstructed masses show very clearly, that this
ml-based method provides conceptually different results from other
methods, that preserve the physical constraints more consistently.
Thus, it might be necessary to investigate which method provides the
most sound results for the particular analysis, one is interested in
and adapt accordingly.

For reconstructed variables, it is always important that they are
well-modelled by the simulations. To ensure such a consistency for
the reconstruction provided by the MultiModalTransformer,
distributions of reconstructed variables such as the top-quark mass
were compared between simulation and data collected by \atlas.
Generally, a very good agreement in the distributions was found.
Highlighting that this tools can be applied to simulation as well as
event data. The distributions show the robustness of the method to be
used on data and in future analyses.

Combined with the thorough evaluation of the method's performance, it
provides a clear motivation to use an approach developed here, or a
similar one, in future research. Lastly, the combination of two
reconstruction tasks into a single model not only resolves
interdependencies but can also cut computational load. The transition
to the HL-LHC, which will increase the event number by up to one
order of magnitude, will bring additional challenges related to the
computation time of analyses. Therefore, tools that combine tasks are
becoming increasingly important as they reduce redundancies.

\section{Outlook}
\label{sec:outlook}
This thesis provided a comprehensive coverage of the strengths and
weaknesses of the methods. Based on the research conducted, new
research goals and questions have emerged.

One very important aspect this thesis highlighted was the
interdependency of the reconstruction tasks in the dilepton channel.
The analysis conducted here showed conclusively that the two tasks
are inherently linked. Thus, it is strongly encouraged that a tool
like $\nu^2$-Flows, which provides a neutrino regression, is extended
to output the assignment. This is motivated by the finding, that a
transformer trained to regress the neutrinos build an internal
understanding about the assignment structure. Hence, a tool should
generally provide both aspect. This would also increase consistency
between the assignment of jets and regression of neutrinos, which
might be lost when interfaced with additional tools.

The next open question is the aforementioned propagating of the
neutrino posterior to derived event variables. It was shown that
averaging over the posterior of event variables instead of computing
event variables from averaged momenta can lead to improvement such as
reduced biases and enhanced resolutions. Furthermore, it might be
useful to incorporate such a posterior of an event variable into a
fit or perhaps the uncertainty. This would provide a more
comprehensive description of the underdetermined neutrinos.

Going one step further, it is also possible to incorporate the
probability distribution of reconstructed variables to the
loss-function itself. Using automatic differentiation to compute the
Jacobian, the probability density of the neutrinos can be transformed
to the probability density of the reconstructed variables. This would
allow one to optimise the model with respect to the variables, that
one is eventually interested in -- i.e. reconstructed variables such
as $m(t\bar{t})$, \chel, or \chan. Similarly, regression these
variables directly could also be investigated. In case where the
interpretability is less important than the separation power of
variable, one might also consider training a discriminator. One
notable example is the search for contributions from non-relativistic
bound state effects. For such an analysis, a discriminating neural
network can be trained to separate events from both states and this
provide one, or possibly multiple, sensitive variables.

Lastly, as it was shown that the tool can provide improvements for
the analyses involving dileptonic decays of top-quark pairs, it can
be expected to be used in future investigations. The most notable one
remains the investigation of bound-state effects at the
\ttbar-production threshold. For a measurement of the cross-section,
a significant improvement in the uncertainty is projected by the use
of the MultiModalTransformer. To facilitate an application in a
measurement by \atlas, the tool was made available in the central
software used for top-related analyses: TopCPToolkit \cite{topcptoolkit}.
Combined with the increased statistics offered by the \LHCruniii in
the near future, very promising progress in the understanding of
these effects can be expected.
